\documentclass[12pt, a4paper]{article}
\usepackage{mathptmx}
\usepackage[T1]{fontenc}
\usepackage[utf8]{inputenc}
\usepackage[brazilian]{babel}
\usepackage[table]{xcolor}
% \usepackage{abntex2}
\usepackage[alf,abnt-etal-cite=1,abnt-etal-list=0, abnt-and-type=e]{Config/abntex2cite}
\usepackage{geometry}
\usepackage{graphicx}
\usepackage{array}
\usepackage{booktabs}
\usepackage{pifont}
\usepackage[final]{pdfpages}
\usepackage[table]{xcolor}
% \usepackage[hidelinks]{hyperref}
% \usepackage{bookmark}

\usepackage{indentfirst}
\usepackage{fontspec}
\setmainfont{Arial}
\usepackage{listings}
\usepackage{setspace}
\usepackage{adjustbox}
\usepackage{varioref}

\newcommand*\rot{\rotatebox{90}}
\newcommand*\V{\ding{51}}

\newcommand{\candidato}[1]{\def\candidatonome{#1}}
\newcommand{\orientador}[1]{\def\orientadornome{#1}}
\newcommand{\coorientador}[1]{\def\coorientadornome{#1}}
\newcommand{\professordisciplina}[1]{\def\professordisciplinanome{#1}}
\newcommand{\titulo}[1]{\def\titulonome{#1}}
\newcommand{\programa}[1]{\def\programanome{#1}}
\newcommand{\mes}[1]{\def\mesnome{#1}}
\newcommand{\ano}[1]{\def\anonome{#1}}
\newcommand{\area}[1]{\def\areanome{#1}}
\newcommand{\linhapesquisa}[1]{\def\linhapesquisanome{#1}}
\newcommand{\citeC}[1]{\citeonline{#1}}
\usepackage{breakcites}
\usepackage{lipsum}
\usepackage{graphicx}
\usepackage{colortbl}
\usepackage{longtable}
\usepackage{amsmath}

\usepackage{fancyhdr}
\usepackage{placeins} % put this in your pre-amble
\usepackage{flafter}  % put this in your pre-amble
\usepackage{enumitem}

\newcolumntype{L}{>{\centering\arraybackslash}m{15cm}}

% Define as dimensoes
\geometry{
	a4paper,
	total={157mm,240mm},
	left=30mm,
	top=20mm,
	right=20mm,
	bottom=40mm
}
\renewcommand{\baselinestretch}{1.5}

\pagestyle{fancy}
\fancyhead[L]{}
\fancyhead[R]{ \includegraphics[width=3cm]{Config/unb.png} }
\renewcommand{\headrulewidth}{0pt}

\begin{document}
%% Informações do candidato
\programa{Pós-Graduação em Geociências Aplicadas e Geodinâmica}
\candidato{Nome do Aluno}
\orientador{Prof. Dr. Nome do Orientador}
\coorientador{Prof. Dr. Nome do Coodientador}
\professordisciplina{Dr. Nome do Professor da Disciplina}
\titulo{Template para a Disciplina Seminário 1 do Instituto de Geociências da Universidade de Brasília}	% Se o título ocupar mais do que duas linhas é preciso modificar o espaçamento no arquivo Config/capa para evitar que a capa ocupe duas páginas
\area{Geoprocessamento e Análise Ambiental}
\linhapesquisa{Avaliação de Dados e Técnicas de Sensoriamento Remoto, Geoprocessamento, Cartografia e Geodésia}
\mes{Agosto}
\ano{2020}

\thispagestyle{empty}
\begin{figure}[h!]
	\flushright
	\includegraphics[width=3cm]{Config/unb.png}
\end{figure}
\begin{center}
	\large
	\MakeUppercase{
		\programanome \\
		Instituto de Geociências \\
		Universidade de Brasília
	}
\end{center}

\vspace{4cm}

\begin{center}
	\Large
	Projeto de Pesquisa de Mestrado \\
	\textbf{\MakeUppercase \titulonome}
\end{center}

\vspace{3.5cm}

\begin{center}
	\large
	\candidatonome \\
	Orientador: \orientadornome \\
	Coorientador: \coorientadornome \\
	\vspace{1.0cm}
	Brasília, \mesnome ~de \anonome
\end{center}
% \keywords{Inteligência Artificial, Mapeamento Sistemático, Extração Automática de Feições}
\clearpage

\thispagestyle{empty}
\begin{figure}[h!]
	\flushright
	\includegraphics[width=3cm]{Config/unb.png}
\end{figure}

\begin{center}
	\large
	\MakeUppercase{
		\programanome \\
		Instituto de Geociências \\
		Universidade de Brasília
	}
\end{center}	

\vspace{1cm}	

\begin{center}
	\MakeUppercase{Projeto de Pesquisa de Mestrado \\
	\titulonome}
\end{center}

\vspace{1cm}	

\begin{center}
	Área de Concentração: \areanome \\
	Linha de Pesquisa: \linhapesquisanome \\
	\candidatonome \\
	Orientador: \orientadornome \\
	Coorientador: \coorientadornome \\
\end{center}

\vspace{1cm}

\hfill \parbox{8.5cm} {\noindent Proposta de Projeto de Mestrado apresentada na Disciplina Seminário I, como pré-requisito do Programa de Pós-Graduação em Geociências Aplicadas e Geodinâmica, do Instituto de Geociências da Universidade de Brasília. \\
Professor: \professordisciplinanome
}
\vfill

\vspace{1cm}

\begin{center}
Brasília - DF

\mes de \ano
\end{center}

\clearpage
\setcounter{page}{1}
\onehalfspacing

% ---
% inserir lista de ilustrações
% ---
% \pdfbookmark[0]{\listfigurename}{lof}

\listoffigures
\fancyfoot{}
% \thispagestyle{empty}
\cleardoublepage
% ---
% ---
% ---
% ---
% inserir lista de tabelas
% ---
% \pdfbookmark[0]{\listtablename}{lot}

\vspace{1cm}
\listoftables
\fancyfoot{}

% \thispagestyle{empty}
\cleardoublepage
% ---
% ---
% inserir lista de abreviaturas e siglas
% ---
% \pdfbookmark[0]{\listadesiglasname}{los}
% \begin{siglas}
%   \item[ABNT] Associação Brasileira de Normas Técnicas
%   \item[abnTeX] ABsurdas Normas para TeX
% \end{siglas}

% \pdfbookmark[0]{\contentsname}{toc}
% \thispagestyle{empty}
\tableofcontents
\fancyfoot{}
\cleardoublepage
\newpage

\fancyfoot[R]{
	\footnotesize
	\thepage
}
\fancyfoot[C]{
	\footnotesize
	PROGRAMA DE PÓS-GRADUAÇÃO EM GEOCIÊNCIAS APLICADAS E GEODINÂMICA \\
	Instituto de Geociências - Campus Universitário Darcy Ribeiro \\
	Brasília, DF - CEP 70910-900
}
\renewcommand{\footrulewidth}{1pt}

% Section I
\section{INTRODUÇÃO}
\label{sec:intro}

Escrever aqui a introdução. Este é um exemplo de citação de frase \cite{Diakogiannis2020}.


Esse é um exemplo de citação corrida: Segundo \citeC{Diakogiannis2020}...    % basic introduction
\section{OBJETIVO}
\label{sec:objetivos}

Essa pesquisa tem por objetivo determinar ...

\subsection{Objetivos específicos}

Como objetivos específicos da pesquisa, pretende-se:
\begin{enumerate}[label=(\alph*)]
	\item Objetivo específico 1;
 	\item Objetivo específico 2;
	\item Objetivo específico 3 ...; 
\end{enumerate}

\section{FUNDAMENTAÇÃO TEÓRICA}
\label{sec:fundamentacao_teorica}

Será elaborada na disciplina Seminário II.
\section{MATERIAIS E MÉTODOS}
\label{sec:mat_e_met}

Apresentação da seção...

\subsection{Subseção exemplo}
\label{sec:subsec_exemplo}

Exemplo de uso de equação: Segundo \citeC{Diakogiannis2020}, \textit{overall accuracy}, \textit{precision}, \textit{recall}, índice $F_1$ e o coeficiente de correlação de Matthews ($MCC$)  são definidos por:

\begin{equation}
OA=\frac{TP+TN}{FP+FN}
\end{equation}

\begin{equation}
precision = \frac{TP}{TP+FP}
\end{equation}

\begin{equation}
recall = \frac{TP}{TP+FN}
\end{equation}

\begin{equation}
F_1 = 2 \times \frac{ precision \times recall }{precision + recall}
\end{equation}

\begin{equation}
MCC = \frac{ TP \times TN - FP \times FN } { \sqrt{(TP+FP)(TP+FN)(TN+FP)(TN+FN)} }
\end{equation}

\noindent
onde $TP$, $FP$, $FN$ e $TN$ são os vertadeiros positivos, falso positivos, falsos negativos e verdadeiros negativos da inferência respectivamente.
\section{RESULTADOS ESPERADOS}
\label{sec:result}

Como resultados da pesquisa em questão, pretende-se:

\begin{enumerate}
    \item Resultado esperado 1;
    \item Resultado esperado 2;
    \item Resultado esperado 3, etc;
    
  \end{enumerate}

\clearpage
\section{CRONOGRAMA}
\label{sec:cronograma}

As atividades da pesquisa foram planejadas de acordo com os dados da Tabela ~\vref{tab:CronogramaPesquisa}. 


\begin{table}[ht]
\centering
\begin{adjustbox}{width=1\textwidth}
\rowcolors{1}{white}{gray!25}
\small\addtolength{\tabcolsep}{-4pt}
\begin{tabular}{|c l |c|c|c|c|c|c|c|c|c|c|c|c|c|c|c|c|c|c|c|c|c|c|c|c|c|c|}
	\hline
    & & \multicolumn{24}{c|}{Meses}\\
    \hline
	\multicolumn{2}{|c|}{Atividades} & 1 & 2 & 3 & 4 & 5 & 6 & 7 & 8 & 9 & 10 & 11 & 12 & 13 & 14 & 15 & 16 & 17 & 18 & 19 & 20 & 21 & 22 & 23 & 24\\
	\hline
	1- & Cursar Disciplinas 				&\V &\V &\V &\V &\V &\V &\V &\V &\V &\V &\V &\V &  &  &  &  &  &  &  &  &  &  &  &  \\
	2- & Revisão Bibliográfica 			&  &  &  & \V & \V & \V & \V &\V &\V  &\V &\V &\V &\V &\V &\V &\V &\V &\V  &  &  &  &  &  &  \\
	3- & Implementação de Algoritmos 	&   &  &  & \V & \V & \V & \V & \V & \V &\V &\V &\V &\V &\V &\V &\V &\V &\V & & & & & & \\
	4- & Treinamento da Rede Neural 		&   &  &  &  & \V & \V & \V & \V & \V & \V &\V &\V &\V &\V &\V &\V &\V &\V &\V &\V & & & & \\
	5- & Análise dos Resultados 			&   &  &  &  &  &  &  &  &  & &   &   &   &  &\V &\V &\V &\V &\V &\V &\V & & & \\
	6- & Elaboração de artigos 			&   &  &  &  &  &  &  &  & \V  & \V & \V  & \V  & \V  & \V & \V & \V & \V &\V &\V &\V &\V &\V &\V & \\
	7- & Elaboração da Dissertação 		&   &  &  &  &  &  &  &  &  &  &  &\V  &\V  &\V  &\V  &\V  &\V   &\V   &\V    &\V &\V &\V &\V & \\
	8- & Entrega da Dissertação 			&   &  &  &  &  &  &  &  &  &  &  &  &  &  &  &  &   &   &   &   &   &   & \V  & \\
	9- & Defesa 			&   &  &  &  &  &  &  &  &  &  &  &  &  &  &  &  &   &   &   &   &   &   &   & \V \\
	\hline
\end{tabular}
\end{adjustbox}
\caption{Cronograma da Pesquisa} 
\label{tab:CronogramaPesquisa}
\end{table} 
\section{VIABILIDADE TÉCNICO-FINANCEIRA}
\label{sec:viabilidade}

Com respeito à viabilidade técnica, ...



% \input{4-cronograma}
\clearpage
\bibliographystyle{Config/abnt-alf}

\section{REFERÊNCIAS BIBLIOGRÁFICAS}
{ %Disable chapter command
	\renewcommand{\section}[2]{}
	\bibliography{Biblio/Biblio}
}
	
\end{document}